\scanpage{55}
\mspage{61}

\deleted{\unsure{B)}} \uline{Notions de} $\unsure{\mathbb D|\underline D}$-\uline{équivalence} et \uline{de} $\mathbb D$-\uline{contractilité}

On va donc reprendre les questions d'acyclicité et de chgt. de base.

\unsure{Déf.} 1 Soit $\mathbb D$ un \inserted{pré}dérivateur, et soit

\diagram{\begin{tikzcd}[ampersand replacement=\&,row sep=small,column sep=small]
f:X \arrow[rr] \arrow[dr,dotted,"p"'] \& {} \& Y \arrow[dl,dotted,"q"] \\
{} \& e \& {}
\end{tikzcd}}

dans $\mathrm{Diag}$. On dit que $f$ est une \uline{$\mathbb D$-équivalence} si le foncteur

\[ F\longrightarrow f^*F \qquad \mathbb D(Y)\longrightarrow\mathbb D(X) \]

\inserted{induit un foncteur} \deleted{\unsure{est}} pleinement fidèle sur la sous-catégorie de $\mathbb D(Y)$ formée des $\mathbb D$-coefficients \uline{constants}, i.e. si pour tt $M,\inserted{N}\in\mathrm{Ob}(A=\mathbb D(e))$ (``coefficients absolus'') on a

\[ \mathrm{Hom}(q^*M,q^*N)\xrightarrow{\sim}\mathrm{Hom}(p^*M,p^*N) \tag*{(1)} \]

Supposons que $p_*$ et $q_*$ existent. On posera, pour un faisceau $F$ sur $X$

\[ H^{\bullet}_{\mathbb D}(X,F)\overset{\text{\unsure{def|déf}}}{=}p_*(F)\qquad
(\text{et de }{=} \ H_{\bullet}^{\mathbb D}(X,\unsure{P|F})=p_!(F)). \]

\margin{\uline{NB} et si $F=p^*M$, $M\in A$, on note simplement $H^\bullet_{\mathbb D}(X,N)$, $H_\bullet^{\mathbb D}(X,N)$}

L'hom (1) s'interprète \inserted{par adj.} comme

\diagram{\begin{tikzcd}[ampersand replacement=\&,column sep=small]
\mathrm{Hom}(M,q_*q^*N) \arrow[r] \arrow[d,equal,"{\deleted{\unsure{\sim}}}"'] \& \mathrm{Hom}(M,p_*p^*N) \arrow[d,equal,"{\deleted{\unsure{\sim}}}"] \\
\mathrm{Hom}(M,H^\bullet_{\mathbb D}(Y,q^*N)) \& \mathrm{Hom}(N,H^\bullet_{\mathbb D}(X,p^*N)) \arrow[l]
\end{tikzcd}}

déduit de l'hom

\[ q_*q^*N\longrightarrow p_*p^*N\qquad\text{i.e. }H^\bullet_{\mathbb D}(Y,N)\longrightarrow H^\bullet_{\mathbb D}(X,N)\tag*{(2)} \]

Donc dire que (1) est iso pour ts $M,N$ signifie que (2) est un iso pour tout $N$. De $=$, si $p_!,q_!$ existent, cela équivaut à dire que

\[ p_{\unsure{!|*}}p^*N.\longrightarrow q_!q^*N)\quad\text{i.e. }H_\bullet^{\mathbb D}(X,N)\longrightarrow H_\bullet^{\mathbb D}(Y,N)\tag*{($2'$)} \]

est \deleted{\unsure{un}} iso \inserted{dans $A$} pour tout $N\in A$
% archive sheet number: 54
